% Chapter 2: Literature Review

\chapter{Literature Review}

\section{Retrieval-Augmented Generation and Its Evaluation}

\subsection{The RAG Paradigm}

Retrieval-Augmented Generation addresses two key limitations of standalone LLMs: hallucinations and the inability to access knowledge beyond their training data. The framework was first explored by \citet{lewis2020retrieval} through the pairing of a dense retriever with a language model. This enabled the outputs to be grounded and derived from external chunks of data fetched at inference time. Also, knowledge storage is distinct from language understanding and can be updated without retraining the model.

This concept has since expanded across several dimensions. \citet{karpukhin2020dense} introduced the concept of dense passage retrieval, which trains encoders for queries and passages to cluster semantically similar content close to each other in embedding space. Later work added re-ranking stages, hybrid sparse-dense retrieval, and iterative retrieval-generation loops. A modern RAG pipeline includes four stages, namely, query understanding, retrieval, re-ranking, and generation.

\subsection{The RAGAS Evaluation Framework}

RAG evaluation is hard for reasons that traditional NLP metrics can't address as of now. BLEU and ROUGE were designed for summarisation and other tasks, where a small set of correct outputs is assumed and used for scoring. RAG outputs are open, and sometimes a single query can have many correct and valid responses. Human evaluation is the obvious strategy, but it is too expensive at the production scale.

RAGAS, first proposed by \citet{es2024ragas}, demonstrated automated evaluation using LLM judges across three metrics. Faithfulness measures whether the response is grounded in the context. Answer relevance and context relevance measure whether the answer is related to the query and whether the context addresses the question. Each metric is scored on a continuous scale. RAGAS is now standard in both academic and business RAG evaluation use cases.

\subsection{Known Limitations of LLM Judges}

The LLM-as-a-Judge concept has a few documented failure modes. Position bias in pairwise comparisons, characterised by \citet{wang2023large}, describes how judges prefer responses based on the order in which they appear. Verbosity bias, explored by \citet{saito2023verbosity}, is the tendency for an LLM judge to give preference to longer responses irrespective of content. Self-enhancement bias, the tendency for a judge to favour responses similar to or from the same family of models, was discovered by \citet{panickssery2024self}. Altogether, these biases suggest that LLM judges are not yet model-agnostic evaluators, and changing the features of the input can skew their scoring assessment.

Calibration also varies across different judges. Some tend to produce continuous scores between 0 and 1, and others just make binary decisions. It was highlighted by \citet{zheng2023judging} that judges from different model families have different scoring techniques and calibration. This complicates cross-judge comparisons and ensemble methods.

\subsection{LLM-as-a-Judge Studies on Contemporary Model Families}

The three judges examined in this thesis, namely GPT-5.4-nano, Gemma-4-26B-A4B-IT, and DeepSeek-V3.2 were all released between Q1 2025 and Q4 2025. Each of them has been examined in some adjacent role as a judge or judge-adjacent classifier, but none has been tested under a corpus-poisoning setting. In this section we summarize prior work.

The closest precedent for using nano-class GPT-5 models as judges is \citet{milad2025gpt5}, who benchmarked twelve configurations of the GPT-5 series, including the GPT-5-nano variants, on a medical-question-answering task. They use a reference-anchored pairwise LLM-as-a-judge framework and report that GPT-5-high outperforms the GPT-5-nano variants on both accuracy and rationale quality. Their work is also a useful example of the verbosity and position bias mitigations needed when deploying nano-class judges in high-stakes domains.

\citet{tan2026empirical} investigate practical LLM-as-a-Judge improvement techniques on RewardBench 2 and evaluate the GPT-5.4 family, including the mini and nano variants used as cost-efficient judges. They systematically test task-specific criteria injection, ensemble scoring, calibration context, and adaptive escalation, reporting that task-specific criteria injection contributes +3.0 percentage points and ensemble scoring contributes +9.8 percentage points to judge accuracy on RewardBench 2, reaching 83.6 percent on their best combined configuration. They also find that score variance in ensembles is a useful signal of judge error for mini-class models but a much weaker signal for nano-class models. Their evaluation is on standard reward modelling, not adversarial settings.

For Gemma-4-26B-A4B-IT, we are not aware of published work characterising it in an absolute-scoring judge role. Gemma 4 was Google's open-weight release at the time of our experiment design and attracted substantial community attention as a strong open alternative to proprietary models in the same size class. We selected it as the open-weight judge in our panel partly for that reason: it represents the kind of model that practitioners are likely to reach for when cost or data-residency constraints rule out closed APIs, but its behaviour as a judge in RAG evaluation has not been documented in the LLM-as-a-Judge literature.

For DeepSeek, we are not aware of published work that uses any DeepSeek model in a judge or judge-adjacent classifier role analogous to the use in this thesis. The model family has been widely evaluated as a generation model on standard benchmarks, but its behaviour as an evaluator under either standard or adversarial conditions has not been characterised in the LLM-as-a-Judge literature.

Two broader works frame the field. \citet{gu2024survey} provide a comprehensive overview of LLM-as-a-Judge techniques, benchmarks, and open problems. More recent work on scoring bias \citep{zhang2026scoring} formally defines rubric order bias, score ID bias, and reference answer score bias, extending the bias taxonomy beyond the position, verbosity, and self-enhancement biases established by \citet{zheng2023judging}.

Collectively, these works characterise two of the three judge families used in this thesis in adjacent contexts: medical QA \citep{milad2025gpt5} and general reward modelling \citep{tan2026empirical}. Gemma-4-26B-A4B-IT and DeepSeek-V3.2 have not been evaluated in any judge role we have found. None of the three judges has been evaluated in a corpus-poisoning setting, which is the gap this thesis addresses.

\section{Corpus Poisoning Attacks on RAG Systems}

\subsection{The PoisonedRAG Framework}

The PoisonedRAG framework, introduced by \citet{zou2025poisonedrag}, represents the first systematic study of corpus poisoning attacks against RAG systems. The core insight of PoisonedRAG is that a successful attack requires a malicious passage to simultaneously satisfy two conditions: the retrieval condition, wherein the passage is ranked highly enough to be retrieved for the target query, and the generation condition, wherein the retrieved passage steers the generator toward the attacker-specified answer.

PoisonedRAG operationalises these two conditions through a two-component passage structure. The first component, designed to satisfy the retrieval condition, consists of the target query itself --- ensuring that the passage will achieve high retrieval similarity through direct lexical and semantic overlap. The second component, designed to satisfy the generation condition, is an LLM-generated passage that supports the target incorrect answer with fluent, plausible-sounding prose. The combination creates passages that are both retrievable and persuasive, achieving attack success rates exceeding 90 percent across multiple benchmarks and generator models.

A particularly important finding from PoisonedRAG is that the black-box variant, which requires no access to the retriever's internal parameters, achieves attack success rates comparable to the white-box variant, which assumes full gradient-based access. This demonstrates that retriever opacity provides limited protection against poisoning attacks, as the attacker can construct effective malicious passages using only heuristic principles of semantic similarity.

\subsection{Universal Adversarial Attacks on LLM Judges}

Short adversarial phrases, when appended to candidate responses, inflate the scores given by LLM judges, regardless of the quality of the response. This was explored by \citet{raina2024llm}. The phrases are found through gradient-based optimisation against a set of target judges and exploit specific features the judge reads. This characteristic, uniform across different queries and responses, suggests it is a structural vulnerability and not instance-specific.

\citet{shi2024judgedeceiver}, through the JudgeDeceiver method, advanced in this direction by developing universal templates that could sway a judge model without access to its internal parameters. This black-box property makes the attacks practically deployable against evaluation systems across the industry, where you don't have access to the model.

\subsection{Distribution Shift and Stylistic Sensitivity}

Robustness of meta-evaluation of LLM safety judges, studied by \citet{eiras2025know}, showed that judges are surprisingly fragile to distribution shifts that humans would find trivial. Even small style changes like bullet points, a storytelling tone, etc., produced large shifts in judge performance; false negative rates went up by values up to 0.24 on the same dataset.

This point goes beyond a safety perspective. If judges are responding to stylistic changes that are irrelevant to content and its quality, they will also react and respond to other forms of context manipulation. RAG evaluation, in particular, is one such case where the passages retrieved can carry patterns or text that can confuse the judge, whether those were placed there adversarially or not.

\section{Defence Mechanisms for RAG Systems}

\subsection{Perplexity-Based Detection}

Perplexity is one of the simplest things you can compute on a passage, and the obvious first thing to try when filtering a corpus. The intuition is that adversarially crafted text, especially text produced by gradient-based optimisation, looks weird statistically and registers as high-perplexity against a clean reference distribution. Set a threshold, drop anything above it, and the malicious content is gone before it reaches the generator or the judge.

But it does not work reliably for these attacks. \citet{zou2025poisonedrag}, in the original PoisonedRAG paper, showed that the perplexity distributions of clean and malicious passages overlap substantially. Choosing a threshold that catches a useful fraction of the poison also drops a lot of clean passages. The black-box variant is even more resistant, because its malicious text is built from two natural-looking pieces: the question itself, which is grammatical, and an LLM-generated passage, which is fluent. Neither piece looks anomalous on a perplexity scan. Later work like the PR-Attack framework explicitly optimises for low perplexity while keeping the attack effective, which closes off the perplexity defence almost entirely. These limitations motivate approaches that combine multiple complementary detection signals rather than relying on a single statistical indicator.

\subsection{Embedding and Clustering Approaches}

TrustRAG \citep{zhou2025trustrag} takes a different angle. The idea is that malicious passages, optimised to be retrieved for a specific target query, end up clustering together in embedding space. Run K-means on the retrieved passages, and the malicious cluster will become separable from the clean retrieval distribution. After clustering, TrustRAG layers cosine similarity and ROUGE checks to confirm which passages in a suspicious cluster are likely poisoned, then uses the LLM's own knowledge to resolve disagreements between the external passages and what the model already believes.

It works when the poison-to-clean ratio is high enough for clusters to form distinctly. In standard settings, where poisoned documents are a minority of retrievals, comparative evaluations report that TrustRAG produces a high false positive rate, flagging many benign passages as poisoned. This is a recurring pattern in the defence literature: aggressive detection can reduce attack success rates but often at the cost of retrieval quality on clean queries. These limitations motivate lightweight approaches that balance detection capability with preservation of clean retrieval quality.

\subsection{RAGuard and Retrieval-Side Filtering}

RAGuard \citep{cheng2025raguard} is closer to the kind of defence studied in this thesis because it works before poisoned text reaches the generator. It first expands the retrieval scope so that the retrieved set contains more clean passages, and then applies chunk-wise perplexity filtering and text-similarity filtering to flag suspicious texts. The useful point here is that the defence is non-parametric, so it does not need a new classifier to be trained for every attack. At the same time, it is still designed mainly for protecting the generator's final answer, whereas our question is whether an upstream filter can protect the evaluator itself when the evaluator is an LLM judge.

\subsection{The CAP Method and Judge-Side Defences}

The Comparative Augmented Prompting (CAP) method works on a different observation: comparative evaluation tasks (which response is better, A or B) hold up better against adversarial pressure than absolute scoring tasks (give this response a score from 0 to 1). CAP wraps a comparative step around the absolute scoring task to inherit some of that robustness. Tested on SummEval and TopicalChat, it lowers attack success rates without hurting clean evaluation quality much.

CAP is a judge-side defence. It modifies what the judge does, not what the judge sees. That puts it on a different layer of the pipeline from the pre-filter we study in this thesis, which sits between the retriever and the judge and modifies the input.

\section{Multi-Signal Anomaly Detection and Classification}

\subsection{DeBERTa and Discriminative Classifiers}

DeBERTa \citep{he2021deberta} is BERT with two changes: disentangled attention, which separates content and position information, and an enhanced mask decoder. DeBERTa-v3 added replaced-token detection during pretraining, and the small variant keeps most of the classification performance with far fewer parameters. We picked DeBERTa-v3-small (140M parameters) for the classifier signal in our pre-filter for three practical reasons. It performs well on short-text classification, which is what we are doing. It trains in about 15 minutes on a free-tier Kaggle T4 GPU for our setup, which made the iteration loop tolerable. It also runs on CPU at inference time, which matters if you want to deploy the filter without specialised hardware.

We also trained a RoBERTa-base classifier \citep{liu2019roberta} as a comparison baseline. RoBERTa is a well-established model with similar parameter counts, and the comparison checks whether the choice of backbone matters much for the poisoned-passage classification task. Results are reported in the ablation analysis of Chapter 4.

\subsection{XGBoost for Feature Aggregation}

XGBoost \citep{chen2016xgboost} is a gradient-boosted tree ensemble. It is the standard tool for combining heterogeneous tabular features into a single decision, and it has become hard to find a Kaggle competition where it does not show up. In our pre-filter, XGBoost takes the five per-passage signal scores as input and outputs a poisoned-or-not probability. We picked it over a linear combination because the signals interact in non-linear ways: a passage with moderate cosine anomaly is suspicious only if the entropy signal also fires, and so on. Trees handle this kind of interaction naturally. They also tolerate missing or zero-valued signals, which matters because the cross-encoder signal is only computed during GPU training and is set to zero at deployment.

\subsection{Cross-Encoders for Relevance}

Cross-encoders score a query-passage pair jointly: the model takes both as input in a single forward pass and produces one relevance score, attending to all token interactions across the pair. Bi-encoders, by contrast, encode the query and the passage independently and compare them with cosine similarity or a dot product. The cross-encoder is more accurate because it can model fine-grained interactions, but it is also slower because passage embeddings cannot be precomputed independently. That trade-off is why cross-encoders are typically used for re-ranking a small candidate set after a fast bi-encoder retrieval, rather than for the initial retrieval itself. Cross-encoders are therefore attractive candidates for relevance estimation in filtering and retrieval-quality pipelines.

\section{Research Gap and Positioning of the Present Work}

The literature reviewed in this chapter reveals a clear gap at the intersection of LLM-as-a-Judge robustness and corpus poisoning. Existing defence mechanisms for Retrieval-Augmented Generation systems primarily focus on the generator or retriever: they aim to prevent malicious passages from influencing the generated answer. Comparatively little attention has been given to the evaluator itself. At the same time, the judge-robustness literature has largely focused on attacks that directly manipulate the judge through adversarial prompts, response formatting, or optimisation-based phrase injection, rather than attacks that indirectly influence evaluation by corrupting the retrieved context.

To our knowledge, no prior work has systematically examined how contemporary LLM judges behave when evaluating responses grounded in poisoned retrieval contexts, nor how upstream filtering interventions alter downstream judge reliability under such conditions. This thesis addresses that gap through a controlled empirical study spanning multiple judges, attack strategies, and filtering approaches. The following chapters describe the experimental methodology and present the resulting analyses.
